\section{三、项目调研与需求分析}

项目立项之初，团队即确立了“先摸清底数、再定义产品”的工作原则，围绕“现状数据调研”与“多方需求调研”两条主线同步开展工作：前者回答“问题有多严重、为什么必须用新的手段解决”，后者回答“系统应该做成什么样、由谁来用”。两条主线的结论彼此印证，共同构成了本项目技术路线与功能设计的现实依据。

\subsection{1.现状数据调研：黑土地退化问题的量化验证}

为精准把握黑土地保护的现实痛点，项目针对黑龙江黑土区开展了专项调研。该区域地处松嫩平原，是我国东北黑土带的核心分布区之一，也是本项目后续试点落地的目标区域，其退化状况具有典型性与代表性。调研以遥感影像解译与地面观测资料相结合的方式展开，将原本依赖定性描述的退化趋势转化为可度量、可比较的量化指标，数据结果如下图所示。

在调研方法上，团队采用“遥感解译为主、地面资料为辅”的复合方式：一方面选取覆盖松嫩平原典型黑土区的多期国产高分二号（GF-2）亚米级光学遥感影像，通过多时相对比提取地表覆盖变化信息，识别土层裸露、侵蚀沟发育与植被覆盖度下降等退化迹象；另一方面收集整理区域内长期积累的地面观测记录与土壤普查资料，对遥感判读结果进行抽样校核，降低单一数据源可能带来的误判风险。在样本构成上，调研兼顾了地形地貌与耕作方式的差异，既包含地势平坦、连片耕作的代表性地块，也覆盖坡耕地与侵蚀沟发育相对集中的敏感区域，使结论既能反映区域整体趋势，也能揭示退化在空间分布上的不均匀性。以下三项指标，即为上述方法下形成的核心研判结论。

\begin{figure}[htbp]
    \centering
    \fbox{\parbox[c][4.5cm][c]{0.8\textwidth}{\centering 【图片空位：image2.png】黑龙江黑土区退化状况专项调研统计数据展示}}
    \caption{黑龙江黑土区退化状况专项调研统计数据}
    \label{fig:img2}
\end{figure}

调研结果显示，黑龙江黑土区退化形势严峻，主要体现在以下三个方面：

\begin{itemize}
    \item \textbf{土层变薄}：黑土层平均减薄厚度达6.83cm，远超自然侵蚀速度。黑土层的形成需要长期腐殖质积累，而流失却可能在数个耕作季内发生，这一减薄幅度直接压缩了作物根系的发育空间与养分供给基础。
    \item \textbf{肥力下降}：土壤有机质含量下降占比达52.5\%，近半数耕地出现明显肥力衰退。有机质降低会同步削弱土壤的保水保肥能力与团粒结构稳定性，往往表现为化肥施用量被动增加而产量提升有限。
    \item \textbf{水土流失}：水土流失问题突出，占比达24.4\%，侵蚀沟扩张对耕地造成直接破坏。侵蚀沟不仅切割地块、阻断农机作业通路，还会将富含有机质的表土持续搬运出境，形成“越冲越薄、越薄越冲”的恶性循环。
\end{itemize}

从成因上看，三项退化指标并非彼此孤立，而是相互耦合、互为因果。长期高强度翻耕使耕层土壤结构被反复扰动，有机质矿化速度加快、团聚体稳定性下降，土壤一旦遭遇降雨便更易被径流带走，这既解释了有机质下降与水土流失在空间上的高度重叠，也说明土层减薄并非单纯的自然过程，而是耕作方式、种植结构与自然侵蚀共同作用的结果。从治理含义上看，这一耦合关系决定了黑土地保护不能停留在单点补救：若只关注肥力指标而忽视侵蚀沟的扩展，施入的有机物料仍会随表土流失而难以积累；若只治理侵蚀沟而不改变耕作方式，新的侵蚀仍会持续发生。因此，治理必须以地块为单元、以时间序列为线索，持续跟踪土层厚度、有机质含量与侵蚀状况的动态变化，才能判断措施是否真正见效并及时调整。这对监测能力提出了新的要求——不仅要“看得清”，还要“看得勤、看得准、看得出趋势”，从而把治理从经验驱动逐步转向数据驱动。

这些量化数据，直观印证了黑土地“变薄、变瘦、变硬”的严峻趋势。需要指出的是，退化问题的紧迫性不仅来自数据本身，更来自监测能力的滞后：传统黑土地监测主要依靠人工实地巡查与逐级上报，存在采样点稀疏、更新周期长、空间覆盖不连续、人力与时间成本高等固有短板，难以在退化发生的早期阶段发现苗头，也无法支撑以地块为单元的精细化治理决策。正是这一“问题在加速、手段在滞后”的错位，为本项目的研发提供了明确的现实依据，也决定了项目必须走“高分辨率遥感 + AI 智能解译 + 快速更新”的技术路线。

\subsection{2.多方需求调研：明确项目落地的实际场景}

为确保项目成果可落地、可推广，项目同步开展了多渠道需求调研，将农户、基层农技人员、土地管理部门以及农业科技企业等不同主体纳入调研范围，力求从“使用者”和“服务者”两个视角交叉验证需求，避免技术方案与真实场景脱节。

两类调研在问题设定上各有侧重：问卷调查面向数量众多的一线使用者，回答的是“需求有多普遍、痛点集中在何处”的问题，其价值在于样本覆盖面广、反馈口径统一，能够从大量分散的意见中提炼出共性诉求；企业走访面向技术供给方与治理实施方，回答的是“现有方案为什么不够用、落地时会受哪些约束”的问题，其价值在于能够接触真实项目中的部署环境、运维条件与成本结构，从而避免方案设计停留在理想条件。前者告诉团队“该做什么”，后者告诉团队“能做成什么”，两条线索相互约束，才共同划定了既先进又可行的设计边界。

\begin{enumerate}
    \item \textbf{问卷调查}：面向农户、基层农技人员、土地管理部门发放问卷，了解一线监测的痛点。调研反馈集中指向三类问题：一是人工巡查效率低，地块分散、路况复杂，单次全覆盖核查耗时长，难以形成常态化监测；二是数据更新慢，从发现异常到逐级汇总形成可用的结论往往周期偏长，错过最佳干预窗口；三是分析成本高，专业遥感影像解译门槛较高，基层单位缺乏相应的技术力量与软件条件。上述痛点共同指向一个明确诉求——亟需一套自动化、可视化的智能监测工具，让不具备遥感专业背景的人员也能看懂结果、用上结果。
    \item \textbf{企业走访}：走访农业科技企业与土地治理相关单位，了解现有监测技术的应用现状与实际约束。调研发现，现有技术方案或依赖重资产的专业平台，部署与运维成本高；或功能单一，仅能提供影像浏览与基础统计，难以支撑从监测到决策的完整链条。据此，项目明确了系统需支持“卫星+无人机”多源数据处理、批量检测、一键生成报告等核心需求，并进一步确立了轻量化、可本地化部署的设计取向，以适配基层单位现有的硬件条件与运维能力。
\end{enumerate}

通过现状数据与多方需求的双重调研，项目明确了“以 AI 遥感监测为核心，打造全流程可视化智能系统”的技术路线。这一路线既回应了退化数据所揭示的紧迫性，也严格对齐了使用者在效率、时效、成本与易用性上的真实约束，确保项目成果真正贴合一线应用场景，解决实际问题。

具体而言，调研结论向技术需求的转化沿三条线索展开：其一，退化现象的空间异质性与地块尺度特征，要求数据源必须具备亚米级空间分辨率与较高的重访频次，这直接指向国产高分二号（GF-2）影像与无人机航拍影像的组合使用，形成“卫星看全局、无人机看重点”的协同观测方式；其二，传统监测“更新慢、覆盖不连续”的短板，要求系统把变化检测而非单期分类作为核心能力，并支撑“周级”的动态更新节奏，这构成了项目选用双时相变化检测算法并以改进 BIT（Bipartite Transformer）为主线的直接原因；其三，基层单位缺乏专业遥感人员与高性能算力的现实约束，要求系统在保证识别精度的前提下尽可能压低部署门槛，由此确定了前后端分离的轻量化架构与可本地化部署的设计原则。经过这一映射，调研发现的问题被逐条转化为可验证的技术指标与功能清单，成为后续模型选型、系统开发与试点验证的共同依据。
